\documentclass[12pt]{article}

\usepackage{sbc-template}

\usepackage{graphicx,url}
\usepackage{pgfplots}

\usepackage[brazil]{babel}
\usepackage[utf8]{inputenc}

\sloppy

\title{Avaliação de Desempenho de Algoritmos Exatos e Heurísticos para o Problema dos k-Centros}

\author{Adriano A. D. dos Santos\inst{1}, Rafael G. Grossi\inst{1}}

\address{Instituto de Ciências Exatas e Informática -- Pontifícia Universidade
    \\ Católica de Minas Gerais (PUC Minas)
    \\ R. Cláudio Manoel, 1162 -- Savassi - 30.140-100 -- Belo Horizonte -- MG -- Brazil
    \email{\{adriano.domingos, rafael.georgetti\}@sga.pucminas.br}
}

\begin{document} 

\maketitle

\section{Implementação}

O grafo é lido de um arquivo em formato de texto pelo \texttt{GraphReader}, que extrai o número de vértices $n$, o número de arestas $m$ e o parâmetro $k$, seguidos de $m$ linhas, sendo que cada uma representa uma aresta ponderada. As arestas são carregadas utilizando a representação \textit{Forward Star}, a qual é utilizada para executar o algoritmo de Floyd-Warshall. Este algoritmo computa as distâncias mínimas entre todos os pares de vértices e produz uma matriz $n \times n$ de caminhos mínimos. Essa matriz é então fornecida como entrada para cada algoritmo solucionador do problema dos $k$-centros (Gonzalez, FastMap, WVA-IG e Exact Solver), que a utilizam para selecionar os centros e calcular o raio da solução.

\subsection{Floyd-Warshall} \label{floyd-warshall}

O algoritmo de Floyd-Warshall computa as distâncias mínimas entre todos os pares de vértices com complexidade de tempo $O(|V|^3)$. Sua implementação segue a formulação padrão: a matriz de distâncias é inicializada com $0$ na diagonal principal e $\infty$ para os demais pares. Em seguida, as arestas do grafo são inseridas, considerando-se a natureza não-direcionada do grafo. Executam-se então três laços aninhados com a relaxação $dist[i][j] = \min(dist[i][j], dist[i][k] + dist[k][j])$. A implementação realiza verificações de continuidade para evitar \textit{overflow} e, ao final, atesta a ausência de ciclos negativos. A matriz de distâncias resultante, contendo os caminhos mínimos entre todos os pares, constitui a entrada para todos os algoritmos de $k$-centros, garantindo que as distâncias satisfaçam a desigualdade triangular.

\subsection{Exact Solver} \label{exact-solver}

O resolvedor exato opera sobre o problema de decisão associado ao $k$-centros: dado um raio $R$, decide-se se existem $k$ centros em $V$ tais que todo vértice esteja a uma distância $\leq R$ de pelo menos um centro. A busca pelo raio mínimo é realizada mediante uma busca binária sobre os valores distintos da matriz de distâncias \cite{kariv1979}, reduzindo o número de problemas de decisão a $O(\log n^2)$.

Para cada raio candidato $R$, constroem-se, para cada vértice $i$, máscaras binárias de 64 bits (\textit{bitsets}) representando o conjunto de vértices $j$ tais que $dist[i][j] \leq R$. O problema de decisão reduz-se a uma instância de \textit{set cover}: existem $k$ centros cujas máscaras, quando combinadas via operação lógica OR (\textit{bitwise}), cobrem todos os $n$ vértices? A exploração exaustiva é realizada por \textit{backtracking}, que tenta todas as combinações de $k$ centros. Para viabilizar a execução em instâncias de porte moderado, quatro técnicas de poda são empregadas:

\begin{enumerate}
    \item \textbf{Eliminação de centros dominados.} Um centro $i$ é dominado por $j$ se o conjunto de cobertura de $i$ está contido no de $j$. Nesse caso, $i$ é redundante, pois $j$ cobre tudo que $i$ cobriria. Essa dominância é detectada por comparação par-a-par das máscaras, removendo-se permanentemente os centros dominados \cite{caruso2003}.

    \item \textbf{Verificação gulosa (\textit{greedy check}).} Antes de iniciar o \textit{backtracking}, aplica-se a heurística gulosa de \cite{chvatal1979} para a cobertura de conjuntos: seleciona-se iterativamente o centro que maximiza o número de vértices recém-cobertos. Se a cobertura completa for obtida em até $k$ passos, o raio $R$ é considerado viável e a recursão é evitada completamente, acelerando a verificação.

    \item \textbf{Ordenação pelo menor número de opções (MRV).} No \textit{backtracking}, seleciona-se o vértice descoberto $u$ com o menor número de centros capazes de cobri-lo (heurística \textit{minimum remaining values}~\cite{haralick1980}). Isso minimiza o fator de ramificação nos níveis iniciais da árvore, podando rapidamente regiões inviáveis do espaço de busca.

    \item \textbf{Ordenação dos candidatos por cobertura.} Dentro da lista de centros que cobrem o vértice $u$, os candidatos são ordenados em ordem decrescente pelo tamanho do seu conjunto de cobertura. Centros que cobrem mais vértices são testados primeiro, aumentando a probabilidade de encontrar uma cobertura viável cedo na busca.
\end{enumerate}

O algoritmo retorna a solução ótima exata para o raio mínimo, correspondente ao menor $R$ para o qual o problema de decisão se prova viável.

\subsection{Gonzalez} \label{gonzalez}

O algoritmo de \cite{gonzalez1985} é uma heurística construtiva que oferece uma 2-aproximação para o problema dos $k$-centros. A execução inicia selecionando arbitrariamente o primeiro vértice (índice 0) como primeiro centro. Em cada iteração subsequente, o vértice mais distante do conjunto de centros já selecionados é escolhido como o próximo centro, repetindo-se esse processo até que $k$ centros sejam selecionados. Sua complexidade é $O(nk)$, sendo extremamente rápido na prática. Seu comportamento é determinístico, produzindo sempre a mesma solução para uma dada instância.

\subsection{FastMap} \label{fast-map}

O algoritmo FastMap~\cite{faloutsos1995} projeta os vértices do grafo em um espaço euclidiano bidimensional, preservando aproximadamente as distâncias originais. A projeção é realizada selecionando-se dois vértices pivôs (os mais distantes entre si) e projetando cada vértice sobre a reta que os conecta. Após a projeção em 2D, o algoritmo aplica a inicialização $k$-means++ para selecionar os $k$ centros no espaço projetado. O $k$-means++ seleciona o primeiro centro aleatoriamente e os demais com probabilidade proporcional ao quadrado da distância ao centro mais próximo. Trata-se de um algoritmo aleatorizado, sendo necessária a utilização de uma semente (\textit{seed}) para garantir reprodutibilidade.

\subsection{WVA-IG} \label{wva-ig}

WVA-IG (\textit{Worst Vertex Adjacent - Iterated Greedy}) é uma metaheurística que combina uma solução inicial obtida pelo algoritmo de Gonzalez com um processo iterativo de refinamento. Em cada iteração, identifica-se o vértice mais distante de seu centro mais próximo ($v^*$, o \textit{pior} vértice), adiciona-se $v^*$ como centro temporário (totalizando $k+1$ centros) e remove-se o centro cuja exclusão resulte no menor aumento de raio. Em seguida, aplica-se uma busca local por trocas (\textit{swap}), testando a substituição de cada centro por um vértice não-centro. Se o algoritmo estagnar por 3 iterações consecutivas sem melhorias, aplica-se uma perturbação aleatória (substituição de centros escolhidos ao acaso). O algoritmo executa até 200 iterações e retorna a melhor solução encontrada.

\section{Metodologia} \label{metodologia}

Os experimentos foram realizados utilizando as 40 instâncias da OR-Library~\cite{beasley1990} para o problema das $p$-medianas, adaptadas para o problema dos $k$-centros. Essas instâncias possuem um número de vértices $n$ variando de 100 a 900, e valores de $k$ proporcionais ao tamanho da instância.

Ressalta-se que as instâncias da OR-Library contêm arestas paralelas com pesos distintos entre um mesmo par de vértices, introduzindo ambiguidade no cálculo das distâncias mínimas. Para determinar a estratégia de tratamento adequada, verificou-se o comportamento esperado pela literatura em relação aos \textit{benchmarks}, constatando-se que os resultados de referência correspondem a uma abordagem linear de leitura do arquivo: para cada par de vértices, o peso da última aresta encontrada é o que prevalece. Os experimentos foram, portanto, conduzidos com essa estratégia.
Cada algoritmo foi executado 5 vezes para cada instância, totalizando 200 execuções por algoritmo. Foi estipulado um tempo limite de 1 hora para a execução de cada rotina, limite este motivado pelo custo computacional inerente ao \textit{Exact Solver}.

\subsection{Métricas de Avaliação}

As métricas utilizadas para a avaliação dos algoritmos foram: (i) \textbf{raio} obtido, definido como a maior distância de um vértice ao centro mais próximo; (ii) \textbf{erro percentual} em relação ao raio de referência; e (iii) \textbf{tempo de execução} do algoritmo em milissegundos, excluindo o tempo de pré-processamento do Floyd-Warshall.

\section{Resultados} \label{resultados}

Os experimentos foram conduzidos conforme a metodologia descrita na Seção \ref{metodologia}. As Tabelas \ref{tab:resultados} e \ref{tab:tempo} apresentam as médias entre as 5 execuções de cada uma das 40 instâncias da OR-Library.

\begin{table}[ht]
\centering
\caption{Raio encontrado para as 40 instâncias da OR-Library. \\ (Para algoritmos não-exatos, os valores correspondem à média das 5 tentativas)}
\label{tab:resultados}
\resizebox{\textwidth}{!}{%
\begin{tabular}{lrrrrrrrr}
\hline
Inst. & $n$ & $k$ & Ótimo & Gonzalez & FastMap & WVA-IG & Exact \\
\hline
pmed1  & 100 & 5   & 127 & 186 (46,46\%) & 175,60 (38,27\%) & 127,20 (0,16\%)          & \textbf{127} \\
pmed2  & 100 & 10  & 98  & 131 (33,67\%) & 152,80 (55,92\%) & \textbf{98} (0,00\%)     & \textbf{98} \\
pmed3  & 100 & 10  & 93  & 154 (65,59\%) & 145,80 (56,77\%) & 96,20 (3,44\%)           & \textbf{93} \\
pmed4  & 100 & 20  & 74  & 112 (51,35\%) & 139,20 (88,11\%) & 76 (2,70\%)              & \textbf{74} \\
pmed5  & 100 & 33  & 48  & 72 (50,00\%)  & 94,40 (96,67\%)  & \textbf{48} (0,00\%)     & \textbf{48} \\
pmed6  & 200 & 5   & 84  & 138 (64,29\%) & 115 (36,90\%)    & 82,60 (-1,67\%)          & \textbf{84} \\
pmed7  & 200 & 10  & 64  & 98 (53,13\%)  & 94,60 (47,81\%)  & 65,20 (1,88\%)           & \textbf{64} \\
pmed8  & 200 & 20  & 55  & 81 (47,27\%)  & 101,80 (85,09\%) & 58,60 (6,55\%)           & \textbf{55} \\
pmed9  & 200 & 40  & 37  & 61 (64,86\%)  & 77,60 (109,73\%) & 40 (8,11\%)              & \textbf{37} \\
pmed10 & 200 & 67  & 20  & 30 (50,00\%)  & 59,20 (196,00\%) & 21,20 (6,00\%)           & \textbf{20} \\
pmed11 & 300 & 5   & 59  & 74 (25,42\%)  & 84,40 (43,05\%)  & 59,40 (0,68\%)           & \textbf{59} \\
pmed12 & 300 & 10  & 51  & 71 (39,22\%)  & 87 (70,59\%)     & 52,20 (2,35\%)           & \textbf{51} \\
pmed13 & 300 & 30  & 35  & 60 (71,43\%)  & 67,20 (92,00\%)  & 38,80 (10,86\%)          & \textbf{35} \\
pmed14 & 300 & 60  & 26  & 40 (53,85\%)  & 61,40 (136,15\%) & 30,60 (17,69\%)          & \textbf{26} \\
pmed15 & 300 & 100 & 18  & 25 (38,89\%)  & 52,80 (193,33\%) & 20 (11,11\%)             & \textbf{18} \\
pmed16 & 400 & 5   & 47  & 84 (78,72\%)  & 64,80 (37,87\%)  & \textbf{45} (-4,26\%)    & \textbf{47} \\
pmed17 & 400 & 10  & 39  & 56 (43,59\%)  & 59 (51,28\%)     & 41,60 (6,66\%)           & \textbf{39} \\
pmed18 & 400 & 40  & 28  & 45 (60,71\%)  & 54,20 (93,57\%)  & \textbf{32,80} (17,15\%) & --- \\
pmed19 & 400 & 80  & 18  & 29 (61,11\%)  & 43 (138,89\%)    & \textbf{23,40} (30,00\%) & --- \\
pmed20 & 400 & 133 & 13  & 19 (46,15\%)  & 38,80 (198,46\%) & 16,60 (27,69\%)          & \textbf{13} \\
pmed21 & 500 & 5   & 40  & 53 (32,50\%)  & 54,60 (36,50\%)  & \textbf{40} (0,00\%)     & \textbf{40} \\
pmed22 & 500 & 10  & 38  & 56 (47,37\%)  & 51,60 (35,79\%)  & 39,60 (4,21\%)           & \textbf{38} \\
pmed23 & 500 & 50  & 22  & 35 (59,09\%)  & 44,80 (103,64\%) & \textbf{27,60} (25,45\%) & --- \\
pmed24 & 500 & 100 & 15  & 23 (53,33\%)  & 33,80 (125,33\%) & \textbf{20} (33,33\%)    & --- \\
pmed25 & 500 & 167 & 11  & 15 (36,36\%)  & 36,40 (230,91\%) & 14,20 (29,09\%)          & \textbf{11} \\
pmed26 & 600 & 5   & 38  & 50 (31,58\%)  & 51,60 (35,79\%)  & 38,60 (1,58\%)           & \textbf{38} \\
pmed27 & 600 & 10  & 32  & 43 (34,38\%)  & 44,20 (38,13\%)  & 34 (6,25\%)              & \textbf{32} \\
pmed28 & 600 & 60  & 18  & 28 (55,56\%)  & 35,80 (98,89\%)  & \textbf{22,40} (24,44\%) & --- \\
pmed29 & 600 & 120 & 13  & 19 (46,15\%)  & 33 (153,84\%)    & \textbf{17} (30,77\%)    & --- \\
pmed30 & 600 & 200 & 9   & 14 (55,56\%)  & 28 (211,11\%)    & \textbf{13,20} (46,66\%) & --- \\
pmed31 & 700 & 5   & 30  & 43 (43,33\%)  & 40,80 (36,00\%)  & \textbf{30} (0,00\%)     & \textbf{30} \\
pmed32 & 700 & 10  & 29  & 45 (55,17\%)  & 65,60 (126,21\%) & 30,60 (5,52\%)           & \textbf{29} \\
pmed33 & 700 & 70  & 15  & 25 (66,67\%)  & 31,20 (108,00\%) & \textbf{20,80} (38,67\%) & --- \\
pmed34 & 700 & 140 & 11  & 17 (54,55\%)  & 29,80 (170,91\%) & \textbf{15,40} (40,00\%) & --- \\
pmed35 & 800 & 5   & 30  & 38 (26,67\%)  & 42 (40,00\%)     & \textbf{30} (0,00\%)     & \textbf{30} \\
pmed36 & 800 & 10  & 27  & 41 (51,85\%)  & 50,80 (88,15\%)  & 29 (7,41\%)              & \textbf{27} \\
pmed37 & 800 & 80  & 15  & 24 (60,00\%)  & 31 (106,67\%)    & \textbf{20,40} (36,00\%) & --- \\
pmed38 & 900 & 5   & 29  & 36 (24,14\%)  & 45,20 (55,86\%)  & \textbf{29} (0,00\%)     & \textbf{29} \\
pmed39 & 900 & 10  & 23  & 35 (52,17\%)  & 33,80 (46,96\%)  & 25 (8,70\%)              & \textbf{23} \\
pmed40 & 900 & 90  & 13  & 21 (61,54\%)  & 25,60 (96,92\%)  & \textbf{18,40} (41,54\%) & --- \\
\hline
\end{tabular}
}
\end{table}

\begin{table}[p]
\centering
\caption{Tempo m\'edio de execu\c{c}\~ao (ms) por inst\^ancia. Para algoritmos n\~ao-determin\'isticos (FastMap e WVA-IG), os valores correspondem \`a m\'edia das 5 tentativas.}
\label{tab:tempo}
\small
\begin{tabular}{lrrrrrrrr}
\hline
Inst. & $n$ & $k$ & Gonzalez & FastMap & WVA-IG & Exact \\
\hline
pmed1  & 100 & 5   & $< 1$ & $< 1$ & 190    & 32 \\
pmed2  & 100 & 10  & $< 1$ & $< 1$ & 741    & 99 \\
pmed3  & 100 & 10  & $< 1$ & $< 1$ & 529    & 131 \\
pmed4  & 100 & 20  & $< 1$ & $< 1$ & 1.537  & 12 \\
pmed5  & 100 & 33  & $< 1$ & $< 1$ & 3.122  & 2 \\
pmed6  & 200 & 5   & $< 1$ & $< 1$ & 771    & 68 \\
pmed7  & 200 & 10  & $< 1$ & $< 1$ & 2.116  & 2.443 \\
pmed8  & 200 & 20  & $< 1$ & $< 1$ & 6.136  & 88.451 \\
pmed9  & 200 & 40  & $< 1$ & $< 1$ & 24.318 & 15.066 \\
pmed10 & 200 & 67  & $< 1$ & 2     & 43.228 & 5 \\
pmed11 & 300 & 5   & $< 1$ & $< 1$ & 1.495  & 47 \\
pmed12 & 300 & 10  & $< 1$ & $< 1$ & 5.793  & 1.870 \\
pmed13 & 300 & 30  & $< 1$ & $< 1$ & 26.687 & 54.150.055 \\
pmed14 & 300 & 60  & $< 1$ & $< 1$ & 91.777 & 496.506 \\
pmed15 & 300 & 100 & $< 1$ & $< 1$ & 162.394 & 2.527 \\
pmed16 & 400 & 5   & $< 1$ & $< 1$ & 2.633  & 113 \\
pmed17 & 400 & 10  & $< 1$ & $< 1$ & 7.952  & 12.701 \\
pmed18 & 400 & 40  & $< 1$ & $< 1$ & 92.161 & --- \\
pmed19 & 400 & 80  & $< 1$ & $< 1$ & 257.374 & --- \\
pmed20 & 400 & 133 & $< 1$ & $< 1$ & 480.379 & 798 \\
pmed21 & 500 & 5   & $< 1$ & $< 1$ & 4.115  & 201 \\
pmed22 & 500 & 10  & $< 1$ & $< 1$ & 11.011 & 70.691 \\
pmed23 & 500 & 50  & $< 1$ & $< 1$ & 194.555 & --- \\
pmed24 & 500 & 100 & $< 1$ & $< 1$ & 555.325 & --- \\
pmed25 & 500 & 167 & $< 1$ & $< 1$ & 1.075.321 & 662.118 \\
pmed26 & 600 & 5   & $< 1$ & $< 1$ & 5.538  & 329 \\
pmed27 & 600 & 10  & $< 1$ & $< 1$ & 22.621 & 414.457 \\
pmed28 & 600 & 60  & $< 1$ & $< 1$ & 367.121 & --- \\
pmed29 & 600 & 120 & $< 1$ & $< 1$ & 1.008.528 & --- \\
pmed30 & 600 & 200 & $< 1$ & $< 1$ & 1.915.908 & --- \\
pmed31 & 700 & 5   & $< 1$ & $< 1$ & 7.665  & 489 \\
pmed32 & 700 & 10  & $< 1$ & $< 1$ & 29.876 & 509.250 \\
pmed33 & 700 & 70  & $< 1$ & $< 1$ & 759.810 & --- \\
pmed34 & 700 & 140 & $< 1$ & $< 1$ & 1.983.417 & --- \\
pmed35 & 800 & 5   & $< 1$ & $< 1$ & 11.775 & 447 \\
pmed36 & 800 & 10  & $< 1$ & $< 1$ & 35.975 & 792.275 \\
pmed37 & 800 & 80  & $< 1$ & $< 1$ & 1.219.323 & --- \\
pmed38 & 900 & 5   & $< 1$ & $< 1$ & 15.543 & 709 \\
pmed39 & 900 & 10  & $< 1$ & $< 1$ & 43.406 & 256.296 \\
pmed40 & 900 & 90  & $< 1$ & $< 1$ & 2.133.335 & --- \\
\hline
\end{tabular}
\end{table}

Os resultados mostram que o algoritmo de Gonzalez obteve um erro médio de 49,84\% em relação ao raio de referência, sendo o mais rápido juntamente com o FastMap (menos de 1 ms). O FastMap apresentou erros maiores, com média de 95,30\%, sendo superado pelo Gonzalez na grande maioria dos casos. O WVA-IG destacou-se como a heurística de melhor qualidade, com erro médio de 13,17\%, frequentemente encontrando soluções ótimas ou muito próximas do ótimo. O resolvedor exato não conseguiu finalizar a execução em todas as 40 instâncias; apenas 29 foram concluídas dentro do tempo limite de 1 hora. Nas instâncias concluídas, o resolvedor exato sempre retornou o valor de referência esperado pela literatura, confirmando a corretude da implementação.

\begin{figure}[htbp]
    \centering
    \caption{Tempo médio de execução em função do número de vértices, para instâncias com $k=5$. Eixo $y$ em escala logarítmica.}
    \begin{tikzpicture}
        \begin{axis}[
            xlabel={Número de Vértices ($|V|$)},
            ylabel={Tempo Médio (ms) -- Escala Logarítmica},
            ymode=log,
            xmin=0, xmax=1000,
            ymin=10, ymax=20000,
            legend pos=south east,
            grid=major,
            grid style={dashed, gray!30},
            width=\textwidth,
            height=8cm,
            mark size=2.5pt
        ]
        
        \addplot[color=blue, mark=square*, thick] coordinates {
            (100, 190)
            (200, 771)
            (300, 1495)
            (400, 2633)
            (500, 4115)
            (600, 5538)
            (700, 7665)
            (800, 11775)
            (900, 15543)
        };
        \addlegendentry{WVA-IG}
        
        \addplot[color=red, mark=triangle*, thick] coordinates {
            (100, 32)
            (200, 68)
            (300, 47)
            (400, 113)
            (500, 201)
            (600, 329)
            (700, 489)
            (800, 447)
            (900, 709)
        };
        \addlegendentry{Exact}
        
        \end{axis}
    \end{tikzpicture}
    \label{fig:tempo-exec-v}
\end{figure}

\begin{figure}[htbp]
    \centering
    \caption{Tempo médio de execução em função do número de centros $k$, para instâncias com $n=300$. Eixo $y$ em escala logarítmica.}
    \begin{tikzpicture}
        \begin{axis}[
            xlabel={Número de Centros ($k$)},
            ylabel={Tempo Médio (ms) -- Escala Logarítmica},
            ymode=log,
            xmin=0, xmax=110,
            ymin=10, ymax=100000000,
            legend pos=south east,
            grid=major,
            grid style={dashed, gray!30},
            width=\textwidth,
            height=8cm,
            mark size=2.5pt
        ]
        
        \addplot[color=blue, mark=square*, thick] coordinates {
            (5, 1495)
            (10, 5793)
            (30, 26687)
            (60, 91777)
            (100, 162394)
        };
        \addlegendentry{WVA-IG}
        
        \addplot[color=red, mark=triangle*, thick] coordinates {
            (5, 47)
            (10, 1870)
            (30, 54150055)
            (60, 496506)
            (100, 2527)
        };
        \addlegendentry{Exact}
        
        \end{axis}
    \end{tikzpicture}
    \label{fig:tempo-exec-k}
\end{figure}

Ressalta-se que o resolvedor exato foi executado com um tempo limite de 1 hora por instância. A instância \texttt{pmed13} ($|V|=300$, $k=30$) foi incluída nos resultados para demonstrar a escalabilidade do custo computacional, tendo levado aproximadamente 15 horas (54.150.055 ms) para encontrar a solução de forma isolada. Diferentemente das demais instâncias não concluídas, que foram abortadas no limite de 1 hora, essa execução foi permitida até o fim para fins de análise. Por esse motivo, seu tempo de execução integra as tabelas e figuras, mas deve ser interpretado como um caso atípico de extrapolação do limite prático adotado no experimento.

A Figura \ref{fig:tempo-exec-v} apresenta o tempo médio de execução em função de $|V|$ para instâncias com $k=5$. O WVA-IG varia de 190 ms ($|V|=100$) a 15,5 s ($|V|=900$), enquanto o resolvedor exato mantém-se abaixo de 1 s em todos os cenários de baixo $k$.

A Figura \ref{fig:tempo-exec-k} ilustra o comportamento em função de $k$ para $n=300$. O WVA-IG apresenta crescimento monotônico de tempo em relação a $k$, variando de 1,5 s ($k=5$) a 162 s ($k=100$). O resolvedor exato, por outro lado, exibe um comportamento não-monotônico: é extremamente rápido para $k$ pequeno (47 ms em $k=5$) e para $k$ elevado (2,5 s em $k=100$), contudo, torna-se drasticamente mais lento em valores intermediários de $k$, atingindo 15 horas em $k=30$ e 8 minutos em $k=60$. Esse fenômeno é uma consequência direta da natureza do problema de \textit{set cover} subjacente: quando $k$ é restrito, a árvore de busca é rasa; quando $k$ é permissivo, a heurística gulosa (\textit{greedy check}) encontra uma cobertura viável quase instantaneamente. O pior caso ocorre em valores intermediários de $k$, onde a profundidade moderada da recursão, aliada à falha da heurística gulosa, força o algoritmo a explorar exaustivamente um espaço de busca colossal. Nos demais cenários abordados, os algoritmos Gonzalez e FastMap executam em menos de 1 ms, confirmando sua adequação para aplicações onde a velocidade absoluta é imperativa.

\section{Conclusão}

Este trabalho apresentou a implementação e a comparação de quatro abordagens para o problema dos $k$-centros em grafos não-direcionados: Gonzalez (2-aproximação determinística), FastMap (projeção 2D associada ao $k$-means++), WVA-IG (metaheurística iterativa) e um resolvedor exato (busca binária ancorada em \textit{set cover}). Os experimentos utilizando as 40 instâncias da OR-Library demonstraram que o WVA-IG oferece o melhor compromisso entre qualidade de solução e tempo de processamento para instâncias de grande porte, frequentemente alcançando resultados ótimos. O algoritmo de Gonzalez, apesar de extremamente veloz, produz soluções com erro médio em torno de 50\%, sendo recomendado apenas para sistemas que exigem latência mínima e toleram aproximações. O FastMap não se mostrou competitivo neste domínio específico, exibindo alta variabilidade e resultados inferiores ao Gonzalez. Por fim, o resolvedor exato demonstra alta eficácia e eficiência para instâncias pequenas ou cenários extremos de $k$ ($|V| \leq 300$), porém sofre com a explosão combinatória em valores intermediários, inviabilizando seu uso irrestrito em instâncias mais complexas.

\bibliographystyle{sbc}
\bibliography{sbc-template}

\end{document}
